\documentclass{article}
\usepackage[UTF8]{ctex}
\usepackage{amsmath,amssymb,amsfonts}
\usepackage{enumitem}
\usepackage{geometry}
\geometry{a4paper, margin=1in}
% === 错题本排版优化（不参与题目标识） ===
\setlength{\parindent}{0pt}
\setlength{\parskip}{0.35em plus 0.10em minus 0.08em}
\setlength{\emergencystretch}{3em}
\setlist{itemsep=0.25em, topsep=0.40em, parsep=0pt}
\allowdisplaybreaks[2]
% === 错题本排版优化结束 ===
\begin{document}
\section*{西南财经大学2026年数学分析真题}
\section*{175．西南财经大学2026年数学分析试卷}

\subsection{题目1}
1．求极限 \(\lim _{n \rightarrow \infty}\left(1-\frac{1}{2^{2}+2}\right)^{\frac{1}{n}}\left(1-\frac{1}{3^{2}+3}\right)^{\frac{1}{n}} \cdots\left(1-\frac{1}{n^{2}+n}\right)^{\frac{1}{n}}\) 。

\subsection{题目2}
2 。求 \(z=\frac{1}{2}\left(x^{4}+y^{4}\right)\) 在条件 \(x+y=a(x, y \geq 0)\) 下的最小值，并在此基础上证明不等式：

\[
\frac{x^{4}+y^{4}}{2} \geq\left(\frac{x+y}{2}\right)^{4}
\]

\subsection{题目3}
3．已知 \(4 \frac{\partial^{2} f}{\partial u^{2}}-\frac{\partial^{2} f}{\partial v^{2}}=1\) ，二元函数 \(g(x, y)=f\left(x^{2}+y^{2}, x y\right)\) ，证明：\(\frac{\partial^{2} g}{\partial x^{2}}-\frac{\partial^{2} g}{\partial y^{2}}=x^{2}-y^{2}\) ．

\subsection{题目4}
4．讨论级数 \(\sum_{n=1}^{\infty} \frac{(-1)^{n}}{n \sqrt[n]{n}}\) 的敛散性，如果收敛，是绝对收敛还是条件收敛？

\subsection{题目5}
5．已知二元函数

\[
f(x, y)= \begin{cases}e^{-(x+y)}, & x y>0 \\ 1, & x y \leq 0\end{cases}
\]

积分区域 \(D=\left\{(x, y)|0 \leq|x|+|y| \leq 1\}\right.\) ，求 \(\iint_{D} f(x, y) \mathrm{d} x \mathrm{~d} y\) ．

\subsection{题目6}
6．已知 \(u_{1}>0, u_{n} e^{u_{n+1}}=e^{u_{n}}-1(n=1,2, \cdots)\) ，证明： \(\lim _{n \rightarrow \infty} u_{n}\) 存在，并求其值．

\begin{itemize}

\subsection{题目7}
\end{itemize}

\subsection{题目7}
原题边界“7．”：已知 \(f(x)\) 二阶连续可微，证明：
\[
\int_{a}^{b} f(x) \mathrm{d} x-\frac{1}{2}(b-a)[f(a)+f(b)]=\frac{1}{2} \int_{a}^{b} f^{\prime \prime}(x)(x-a)(x-b) \mathrm{d} x
\]

\subsection{题目8}

\subsection{题目8}
原题边界“8．”：设函数 \(f\) 在 \([0,+\infty)\) 上连续可导，\(f(0)>0\) 且 \(f^{\prime}(x) \geq 0, x \in[0,+\infty)\) ，证明：若 \(\int_{0}^{+\infty} \frac{1}{f(x)+f^{\prime}(x)} \mathrm{d} x\)收敛，则反常积分 \(\int_{0}^{+\infty} \frac{1}{f(x)} \mathrm{d} x\) 也收敛。
\end{document}
